\chapter{Divisorial currents on compact modifications}
\label{chap:divisorial-models}

The numerical trace data of \cref{chap:trace-data} can be assembled into currents on every compact modification. Gauduchon metrics provide exactly the required summability. The construction uses linear mass and does not require a bound for any top-degree Monge--Amp\`ere measure. On each fixed model, the resulting positive closed divisorial current also determines an ordinary potential type with a smooth residual term. This describes the divisorial part on that model; exact realization on all further models is a stronger requirement. A scalar Gauduchon equation gives continuous smooth representatives and proves continuity of closed-current pullback. The difference between the trace currents on two models is then an effective exceptional divisor, whose logarithmic potential describes the change of these ordinary types.

Let $X$ be a connected compact complex manifold, let $Y\subset X$ be a connected smooth compact complex submanifold of dimension $r\geq1$, and fix a smooth real $(1,1)$-form $\theta$ and a Hermitian form $\omega$ on $X$. Let $u\in\PSH(X,\theta)$ satisfy $\nu(u,Y)=\inf_{y\in Y}\nu(u,y)=0$. Use the multiplier ideal $\mathcal I(mu)$ for the weight $\mathrm{e}^{-mu}$ and the trace numbers $\tau_{Y,u}(E)$ of \eqref{eq:trace-profile}. For a smooth compact modification $\pi:Z\to Y$, put
\[
 \theta_Z=\pi^*(\theta|_Y),\qquad
 \omega_Z^0=\pi^*(\omega|_Y),\qquad
 b_m(E)=\operatorname{ord}_E\bigl(\mathcal I(mu)\mathcal O_Y\bigr).
\]
The form $\omega_Z^0$ is only semipositive; a separate Hermitian metric on $Z$ will be used when necessary.

\section{Removing divisorial mass}

For a prime divisor $E$ on a complex manifold, let $[E]$ be its integration current and $\nu(w,E)$ the generic Lelong number of a quasi-plurisubharmonic function $w$ along $E$.

We use the following precise local consequence of Siu's decomposition, as recalled in \cite{SGPT} and the standard current theory of \cite{DemAM11}. A closed positive $(1,1)$-current has, on a prime hypersurface $E$, coefficient-measure restriction equal to its generic Lelong number times $[E]$. In particular the remaining current after removal of any finite set of its divisorial components is positive. Distinct prime hypersurfaces do not contribute overlapping mass: their intersections have codimension at least two. These are statements about the local closed current.

\Needspace{8\baselineskip}
\begin{lemma}\label{lem:nonclosed-divisorial-removal}
Let $M$ be a complex manifold, let $\beta$ be a smooth real $(1,1)$-form on $M$, and let $w\in\PSH(M,\beta)$, and let $E_1,\ldots,E_N$ be distinct prime divisors. Then
\[
 \beta+\ddc w-\sum_{j=1}^N\nu(w,E_j)[E_j]\geq0.
\]
The form $\beta$ need not be closed.
\end{lemma}
\begin{proof}
Work on a coordinate ball and choose a smooth $\rho$ with $\ddc\rho\geq\beta$. The current
\[
 T=\ddc(w+\rho)=(\beta+\ddc w)+(\ddc\rho-\beta)
\]
is closed and positive. A smooth form, viewed through its coefficient measures, assigns zero mass to every proper analytic set. Thus $T$ and $\beta+\ddc w$ have the same restriction to $\bigcup_jE_j$. The stated local Siu fact identifies this restriction as $\sum_j\nu(w,E_j)[E_j]$. Subtracting it from the positive current $\beta+\ddc w$ leaves its coefficient-measure restriction to the complement, which is positive. This proves the inequality locally, hence globally. The argument does not identify $\beta+\ddc w$ itself as a closed current.
\end{proof}

\section{A Gauduchon bound on every model}\label{sec:divisorial-model-bound}

Let $X$ be a connected compact complex manifold, let $Y\subset X$ be a connected smooth compact complex submanifold of dimension $r\geq1$, and fix a smooth real $(1,1)$-form $\theta$ and a Hermitian form $\omega$ on $X$. Let $u\in\PSH(X,\theta)$ satisfy $\nu(u,Y)=\inf_{y\in Y}\nu(u,y)=0$. Use the multiplier ideal $\mathcal I(mu)$ for the weight $\mathrm{e}^{-mu}$ and the trace numbers $\tau_{Y,u}(E)$ of \eqref{eq:trace-profile}. Fix a smooth compact modification $\pi:Z\to Y$, and put $\theta_Z=\pi^*(\theta|_Y)$ and $\omega_Z^0=\pi^*(\omega|_Y)$. For a prime divisor $E$ on any such model, let $b_m(E)=\operatorname{ord}_E(\mathcal I(mu)\mathcal O_Y)$ be the generic order of the pulled-back image ideal. Write $[E]$ for its integration current. Gauduchon metrics $\gamma$ on $Z$ satisfy $\ddc\gamma^{r-1}=0$.

For a positive integer $m$, with $E$ ranging over prime divisors on $Z$, define the finite divisorial current
\[
 D_{m,Z}=\frac1m\sum_E b_m(E)[E].
\]
The sum is finite because the image ideal on the smooth compact manifold $Z$ is coherent and nonzero. Its codimension-one zero components form a finite family. Define initially through finite subsums
\[
 S_Z(u;Y)=\sum_E\tau_{Y,u}(E)[E].
\]

\begin{theorem}\label{thm:trace-divisorial-current-budget}
The last sum defines a closed positive $(1,1)$-current with finite mass on $Z$. For every Gauduchon metric $\gamma$ on $Z$,
\begin{equation}\label{eq:trace-divisorial-linear-budget}
 \int_Z S_Z(u;Y)\wedge\gamma^{r-1}
 \leq\int_Z\theta_Z\wedge\gamma^{r-1}.
\end{equation}
For $r=1$ take $\gamma^{r-1}=1$. Moreover
\[
 0\leq D_{m,Z}\leq S_Z(u;Y),\qquad
 \int_Z\bigl(S_Z(u;Y)-D_{m,Z}\bigr)\wedge\gamma^{r-1}
 \longrightarrow0.
\]
In particular $D_{m,Z}$ converges weakly as a current to $S_Z(u;Y)$. The dyadic sequence $D_{2^j,Z}$ increases.
\end{theorem}
\begin{proof}
Let $u_m$ with curvature errors $\epsilon_m$ be the analytic approximants of \cref{cor:nonclosed-analytic-approximation} and set $w_m=(u_m|_Y)\circ\pi$. They are genuine, and
\[
 T_m=\theta_Z+\epsilon_m\omega_Z^0+\ddc w_m\geq0,
 \qquad \nu(w_m,E)=b_m(E)/m.
\]
\Needspace{6\baselineskip}
The removal lemma gives $D_{m,Z}\leq T_m$. Current Stokes and the Gauduchon identity therefore give
\[
 \int_ZD_{m,Z}\wedge\gamma^{r-1}
 \leq\int_Z\theta_Z\wedge\gamma^{r-1}
       +\epsilon_m\int_Z\omega_Z^0\wedge\gamma^{r-1}.
\]
Pass to the limit in the coefficients of any fixed finite set of prime divisors, using \cref{thm:intrinsic-trace-profile}. Then take the supremum over finite sets. This proves \eqref{eq:trace-divisorial-linear-budget} as a finite-subsum bound.

Every prime divisor has strictly positive $\gamma$-area. Hence there are at most countably many positive coefficients: for each positive integer $k$, at most finitely many terms can have weighted mass at least $1/k$, and their union includes every positive term. Enumerate these divisors. Their finite partial sums are increasing positive closed currents with a uniform trace-mass bound. They therefore converge, in coefficient measures and weakly as currents, to a positive current of finite mass. Distributional differentiation is continuous on currents, so the limit is closed. This constructs $S_Z(u;Y)$ and proves the integral inequality.

The supremum formula for $\tau_{Y,u}(E)$ gives $b_m(E)/m\leq\tau_{Y,u}(E)$ for every $m$ and $E$. The full integer limit gives coefficientwise convergence. Apply dominated convergence to the countable series with summable weights $\tau_{Y,u}(E)\int_E\gamma^{r-1}$. It yields the asserted vanishing trace mass of $S_Z-D_{m,Z}$. A fixed smooth test form is bounded in absolute value on positive $(1,1)$-currents by a constant times this trace mass, proving weak convergence. Finally $b_{2m}(E)\geq2b_m(E)$ gives $D_{2m,Z}\geq D_{m,Z}$.
\end{proof}

The estimate holds for every fixed compact model. It is not a uniform bound over all modifications and all choices of their metrics. Nor does this current convergence assert convergence of nonlinear products of the $D_{m,Z}$.

\Needspace{20\baselineskip}
\section{Compatibility between models}

Let $X$ be a connected compact complex manifold, let $Y\subset X$ be a connected smooth compact complex submanifold of dimension $r\geq1$, and fix a smooth real $(1,1)$-form $\theta$ and a Hermitian form $\omega$ on $X$. Let $u\in\PSH(X,\theta)$ satisfy $\nu(u,Y)=\inf_{y\in Y}\nu(u,y)=0$. Use the multiplier ideal $\mathcal I(mu)$ for the weight $\mathrm{e}^{-mu}$ and the trace numbers $\tau_{Y,u}(E)$ of \eqref{eq:trace-profile}. Fix a smooth compact modification $\pi:Z\to Y$, and put $\theta_Z=\pi^*(\theta|_Y)$ and $\omega_Z^0=\pi^*(\omega|_Y)$. For a prime divisor $E$ on any such model, let $b_m(E)=\operatorname{ord}_E(\mathcal I(mu)\mathcal O_Y)$ be the generic order of the pulled-back image ideal. Write $[E]$ for its integration current. Use the finite divisorial approximants $D_{m,Z}$ and closed positive trace current $S_Z(u;Y)$ defined in \cref{sec:divisorial-model-bound}.

\begin{theorem}\label{thm:trace-divisorial-model-compatibility}
Let $g:Z'\to Z$ be a smooth compact modification, and view both manifolds as models over $Y$. Then
\[
 g_*S_{Z'}(u;Y)=S_Z(u;Y).
\]
The assignment is additive and positively homogeneous in $u$. For every finite set $\mathcal F$ of prime divisors on $Z$, the finite effective real divisor
\[
 A_{\mathcal F}=\sum_{E\in\mathcal F}\tau_{Y,u}(E)E
\]
satisfies
\[
 [g^*A_{\mathcal F}]\leq S_{Z'}(u;Y).
\]
\end{theorem}
\begin{proof}
If $E'$ is the strict transform of a prime divisor $E\subset Z$, the map is biholomorphic near their generic points. They define the same vanishing order on the image ideals from $Y$, and hence $\tau_{Y,u}(E')=\tau_{Y,u}(E)$. Also $g_*[E']=[E]$.

For an exceptional prime divisor $F$, its image has codimension at least two. The pullback of any smooth $(r-1,r-1)$ test form restricts to zero on its regular locus, because the map there has complex rank at most $r-2$. Thus $g_*[F]=0$. In dimension one a modification between smooth compact curves is an isomorphism, so there are no such exceptional divisors. Pushforward is continuous on currents for this proper map. Applying it to increasing finite sums, or equivalently testing their summable positive measures, proves the displayed compatibility. Additivity and homogeneity follow coefficientwise from \cref{thm:trace-data-additivity} and \cref{prop:trace-profile-calculus}.

For the final inequality, let $\mathfrak a_{m,Z}=\mathcal I(mu)\mathcal O_Z$ be the image ideal on $Z$, for a positive integer $m$. Its divisorial vanishing gives
\[
 \mathfrak a_{m,Z}\subseteq
 \mathcal O_Z\left(-\sum_{E\in\mathcal F}b_m(E)E\right).
\]
Indeed on a smooth manifold a holomorphic function with the indicated generic orders is divisible by the corresponding local Cartier equations, with multiplicity. Pull back this inclusion and take the order at any prime divisor $F\subset Z'$. It gives
\[
 b_m(F)\geq\sum_{E\in\mathcal F}
                 b_m(E)\operatorname{ord}_F(g^*E).
\]
Divide by $m$ and pass to the limit; the sum is finite. Thus every coefficient of $g^*A_{\mathcal F}$ is at most $\tau_{Y,u}(F)$. Both sides are divisorial positive currents, which proves the asserted current inequality.
\end{proof}

The argument so far concerns finite divisors. \Cref{sec:continuous-bc-representatives} will prove the required continuity for positive closed $(1,1)$-currents, and \cref{thm:trace-exceptional-difference} will extend the inequality to the whole countable trace current. No such assertion about general currents or nonlinear products is used here.

\Needspace{19\baselineskip}
\section{A majorant in the original restricted background}

Let $X$ be a connected compact complex manifold, let $Y\subset X$ be a connected smooth compact complex submanifold of dimension $r\geq1$, and fix a smooth real $(1,1)$-form $\theta$ and a Hermitian form $\omega$ on $X$. Let $u\in\PSH(X,\theta)$ satisfy $\nu(u,Y)=\inf_{y\in Y}\nu(u,y)=0$. Use the multiplier ideal $\mathcal I(mu)$ for the weight $\mathrm{e}^{-mu}$ and the trace numbers $\tau_{Y,u}(E)$ of \eqref{eq:trace-profile}. Fix a smooth compact modification $\pi:Z\to Y$, and put $\theta_Z=\pi^*(\theta|_Y)$ and $\omega_Z^0=\pi^*(\omega|_Y)$. For a prime divisor $E$ on any such model, let $b_m(E)=\operatorname{ord}_E(\mathcal I(mu)\mathcal O_Y)$ be the generic order of the pulled-back image ideal. Write $[E]$ for its integration current. Use the finite divisorial approximants $D_{m,Z}$ and closed positive trace current $S_Z(u;Y)$ defined in \cref{sec:divisorial-model-bound}. Retain the analytic approximants $u_m$ and errors $\epsilon_m$ of \cref{cor:nonclosed-analytic-approximation}, using $\omega$ for their curvature error.

The preceding currents need not be the complete divisorial data of a single potential. Nevertheless there is one potential that dominates their coefficients simultaneously on all compact models.

\begin{proposition}\label{prop:original-background-trace-majorant}
There exists $v\in\PSH(Y,\theta|_Y)$ such that for every smooth compact modification $\pi:Z\to Y$,
\[
 \nu(v\circ\pi,E)\geq\tau_{Y,u}(E)
 \quad\hbox{for every prime }E\subset Z,
\]
and
\[
 \pi^*(\theta|_Y)+\ddc(v\circ\pi)\geq S_Z(u;Y).
\]
Neither uniqueness nor equality of the individual Lelong numbers is asserted.
\end{proposition}
\begin{proof}
Normalize the genuine restrictions $u_m|_Y$ to have supremum zero. They lie in a fixed Hermitian plurisubharmonic family, because their curvature lower bounds are $-(\theta+\epsilon_m\omega)|_Y$. Normalized compactness gives a subsequence converging in $L^1(Y)$ to a genuine quasi-plurisubharmonic function $v$. Passing to the distributional curvature inequality and using $\epsilon_m\to0$ gives $\theta|_Y+\ddc v\geq0$. Thus the limiting background is the original restricted form.

Fix any compact modification $\pi$. The pullback $L^1$ lemma proved in \cref{lem:modification-lone-pullback} applies to this same subsequence in its fixed smooth lower bound. It gives $L^1$ convergence of the normalized pullbacks to $v\circ\pi$. Upper semicontinuity of generic Lelong numbers under such convergence therefore gives
\[
 \nu(v\circ\pi,E)
 \geq\limsup_m\nu((u_m|_Y)\circ\pi,E)
 =\tau_{Y,u}(E).
\]
Normalization by constants has not affected these numbers. The argument works for each fixed model and each divisor using the already chosen subsequence; no new subsequence is selected for any of them.

Apply \cref{lem:nonclosed-divisorial-removal} to the positive current $\theta_Z+\ddc(v\circ\pi)$. Its divisorial coefficients dominate those of $S_Z$. Removal of every finite partial sum of $S_Z$ preserves positivity. Passing to the current limit from \cref{thm:trace-divisorial-current-budget} proves the final inequality.
\end{proof}

The potential $v$ is a majorant of the numerical orders. Its additional divisorial singularities, if present, have not been excluded. Thus this result does not settle the general exact-trace question of \cref{chap:trace-data}. The analytic and curve-fibration cases there use separately constructed exact representatives. It does explain why all the Gauduchon inequalities in \eqref{eq:trace-divisorial-linear-budget} are compatible with one genuine potential in the original background.

\Needspace{12\baselineskip}
\section{Bott--Chern classes and closed currents}\label{sec:bc-closed-currents}

Bott--Chern cohomology separates the smooth background of a closed current from its global potential. Subtracting the trace current from a closed background then gives a residual class on each model. Pushforward compatibility organizes these classes into a b-divisor class, as in \cite[section \emph{The notion of b-divisors}]{SGPT}.

\index{Bott--Chern class!definition}For a complex manifold $Z$, let $\Lambda^{1,1}_{\mathbb R}$ denote the bundle of real $(1,1)$-forms and define
\[
 H^{1,1}_{\mathrm{BC}}(Z,\mathbb R)
 =\frac{\{\beta\in C^\infty(Z,\Lambda^{1,1}_{\mathbb R}):\mathrm{d}\beta=0\}}
        {\{\ddc f:f\in C^\infty(Z,\mathbb R)\}}.
\]
\Needspace{8\baselineskip}
\begin{lemma}\label{lem:bc-one-one-current-description}
Replacing smooth forms and smooth functions in this quotient by real currents and real distributions gives the same space. In particular every closed positive $(1,1)$-current $T$ on $Z$ can be written
\[
 T=\beta+\ddc w
\]
with $\beta$ smooth and closed and $w$ a global quasi-plurisubharmonic function.
\end{lemma}
\begin{proof}
The decomposition in \cref{lem:closed-current-global-potential} proves surjectivity of the smooth-to-current map and gives a quasi-plurisubharmonic potential for positive currents. Its smooth-difference assertion also proves injectivity: a distribution whose $\ddc$ is smooth is smooth, so two smooth representatives identified in the current quotient are already identified in the smooth quotient.
\end{proof}

\index{Pseudoeffective class}For global pairings, take $Z$ connected and compact. The class of a closed current $T$ is written $\{T\}_{\mathrm{BC}}$. It is called pseudoeffective when it has a positive closed representative. Pullback of classes uses smooth representatives. Pushforward under a proper modification uses currents; it is well-defined because pushforward commutes with $\ddc$. The lemma identifies its result with a smooth Bott--Chern class. For a modification $g$, we have $g_*g^*\beta=\beta$ as currents for every smooth form $\beta$ of the indicated degree, by degree-one change of variables.

\Needspace{9\baselineskip}
A Gauduchon metric $\gamma$ defines a linear functional on this cohomology:
\[
 \ell_\gamma(\{T\}_{\mathrm{BC}})=\int_ZT\wedge\gamma^{\dim Z-1}.
\]
Current Stokes makes it independent of the representative. For a positive representative it is nonnegative, and equality forces that representative to vanish, by positivity and the strict positivity of the metric. In particular the pseudoeffective cone contains no nonzero line: if $T_1,T_2\geq0$ represent opposite classes, their sum has zero Gauduchon pairing, hence $T_1=T_2=0$.

\Needspace{12\baselineskip}
Use the Hopf manifold $X_n$, $n\geq2$, and its Gauduchon metric $\omega$ from \cref{prop:hopf-bergman-exact}, with covering coordinates $z\in\mathbb C^n\setminus\{0\}$. The form $\alpha=\ddc\log|z|^2$ is exact as a de Rham form. Indeed $\dc\log|z|^2$ descends to the quotient because the logarithm changes by a constant under the contraction. Nevertheless
\[
 \ell_\omega(\{\alpha\}_{\mathrm{BC}})
 =\int_{X_n}\alpha\wedge\omega^{n-1}=2(n-1)>0.
\]
Thus its Bott--Chern class is nonzero. Gauduchon integration detects closed-curvature information that the de Rham class loses in this example.

Return now to the smooth compact modifications of the original manifold $Y$.

\Needspace{11\baselineskip}
\begin{definition}\label{def:bott-chern-bdivisor-class}\index{Bott--Chern class!b-divisor}
A \emph{Bott--Chern b-divisor class} over $Y$ is a family
\[
 \mathbb A=(\mathbb A_Z)_{\pi:Z\to Y},\qquad
 \mathbb A_Z\in H^{1,1}_{\mathrm{BC}}(Z,\mathbb R),
\]
over smooth compact modifications, compatible with pushforward. It is pseudoeffective if every $\mathbb A_Z$ is pseudoeffective. It is \emph{Cartier} if there are one model $Z$ and one class $a_Z$ such that its class on every model $Z'$ dominating $Z$ is the pullback of $a_Z$.
\end{definition}

When the models are K\"ahler, the standard $\partial\bar\partial$ lemma identifies this cohomology with the real $(1,1)$-cohomology used in \cite[section \emph{The notion of b-divisors}]{SGPT}. The definition then agrees with the source definition.

For the trace data fixed at the beginning of the chapter, the common majorant supplies positive representatives of the residual classes.

\begin{proposition}\label{prop:trace-bott-chern-bclass}\index{Bott--Chern class!residual trace}
If the original ambient background $\theta$ is closed, the classes
\[
 \mathbb A_Z(u;Y,\theta)
 =\{\pi^*(\theta|_Y)-S_Z(u;Y)\}_{\mathrm{BC}}
\]
form a pseudoeffective Bott--Chern b-divisor class. They depend only on the ambient multiplier-ideal class of $u$ and the Bott--Chern class of $\theta|_Y$.
\end{proposition}
\begin{proof}
By \cref{prop:original-background-trace-majorant}, there is one $v$ on $Y$ such that
\[
 R_Z=\pi^*(\theta|_Y)+\ddc(v\circ\pi)-S_Z(u;Y)\geq0
\]
on every model. Since $\theta$ and $S_Z$ are closed, $R_Z$ is closed. Its class is exactly the displayed one, proving pseudoeffectivity. Compatibility follows from \cref{thm:trace-divisorial-model-compatibility} and $g_*g^*(\theta_Z)=\theta_Z$. The trace data is unchanged by ambient multiplier-ideal equivalence, and replacing $\theta|_Y$ by a smooth $\ddc$ perturbation changes no class. These facts prove independence.
\end{proof}

\Needspace{7\baselineskip}
If $\theta$ is not closed, the same expression $R_Z$ is still a positive current, but
\[
 \mathrm{d}R_Z=\pi^*(\mathrm{d}\theta|_Y).
\]
It therefore need not define a class in ordinary Bott--Chern cohomology. The larger classes with a specified non-closed derivative will be treated in \cref{chap:nonclosed-bclasses}.

\Needspace{12\baselineskip}
\section{The b-divisor class of a closed positive current}

For a closed positive current, the same construction uses its own local potentials in place of restricted ambient ideals. The trace coefficients are then its divisorial Lelong numbers, and divisorial removal gives the positive residual currents directly.

Let $T$ be a closed positive $(1,1)$-current on a connected compact complex manifold $Y$, and write $T=\beta+\ddc w$ as above. On a compact modification $\pi:Z\to Y$, local psh potentials define the positive pullback $\pi^*T$: dominance prevents their compositions from being identically $-\infty$. With $E$ ranging over prime divisors on $Z$, set
\[
 N_Z(T)=\sum_E\nu(w\circ\pi,E)[E],\qquad
 \operatorname{Reg}_Z(T)=\pi^*T-N_Z(T).
\]
Divisorial removal makes the second current positive and closed. The following lemma identifies the coefficients with the trace construction.

\begin{lemma}\label{lem:self-trace-is-siu}
For every quasi-plurisubharmonic function $w$ on $Y$ and every prime divisor $E$ on a smooth compact modification $\pi:Z\to Y$, the self-trace in \eqref{eq:trace-profile} satisfies
\[
 \tau_{Y,w}(E)=\nu(w\circ\pi,E).
\]
In particular, for the potential of the closed positive current $T$, one has $N_Z(T)=S_Z(w;Y)$.
\end{lemma}
\begin{proof}
A genuine quasi-plurisubharmonic function has zero Lelong number almost everywhere, so its self-trace satisfies the required generic hypothesis. For a fixed divisor $E$ on a compact modification, apply \cref{lem:bergman-divisorial-error} on a coordinate neighborhood of its center. With $\nu_E=\nu(w\circ\pi,E)$ and $A_E$ the log discrepancy, it gives
\[
 \max\{\nu_E-A_E/m,0\}
 \leq\frac1m\operatorname{ord}_E\mathcal I(mw)
 \leq\nu_E.
\]
The integer limit defining $\tau_{Y,w}(E)$ is therefore $\nu_E$. Equality of all coefficients identifies the two divisorial currents.
\end{proof}

\Needspace{13\baselineskip}
\begin{theorem}\label{thm:closed-current-bclass}
The family
\[
 \mathbb D(T)_Z=\{\operatorname{Reg}_Z(T)\}_{\mathrm{BC}}
\]
is a pseudoeffective Bott--Chern b-divisor class. For closed positive currents $T_1,T_2$ and $a\geq0$,
\[
 \mathbb D(T_1+T_2)=\mathbb D(T_1)+\mathbb D(T_2),
 \qquad \mathbb D(aT_1)=a\mathbb D(T_1).
\]
On a K\"ahler manifold, with K\"ahler models, it agrees with the b-divisor class of the same current in \cite[subsection \emph{The b-divisors of currents}]{SGPT}.
\end{theorem}
\begin{proof}
The preceding lemma identifies $N_Z(T)$ with the already constructed compatible divisorial currents. For $g:Z'\to Z$, we have $g_*N_{Z'}(T)=N_Z(T)$ and $g_*g^*(\pi^*T)=\pi^*T$. For the latter identity, write the current as a smooth form plus $\ddc$ of a potential; degree-one change of variables applies both to the smooth form and to the locally integrable pulled-back potential, and pushforward commutes with $\ddc$. Therefore
\[
 g_*\operatorname{Reg}_{Z'}(T)=\operatorname{Reg}_Z(T).
\]
Taking classes proves compatibility, and positivity proves pseudoeffectivity. Generic Lelong numbers add under addition of positive currents and scale by nonnegative constants. Hence their divisorial parts, and then their regular parts, have the same properties. This proves the two identities. In the K\"ahler case both the regular-current construction and the identified cohomology classes are exactly those in the source definition.
\end{proof}

\section{Nef Cartier realizations}\label{sec:nef-cartier-realizations}

If $T$ has analytic singularities, a resolution removes them in a finite divisor and leaves a current with bounded potentials. This remainder acquires no new divisorial part on higher models, so its class determines the entire family there. Smooth regularization of its bounded potential makes the realizing class nef.

\index{Nef class!definition}A class $a\in H^{1,1}_{\mathrm{BC}}(Z,\mathbb R)$ is \emph{nef} if, for one Hermitian metric $\gamma$, it has a smooth representative $\beta_\epsilon\geq-\epsilon\gamma$ for every $\epsilon>0$. Comparison of two fixed Hermitian metrics shows that this definition is independent of the metric. A Cartier b-divisor class is called nef here if it has a realization by a nef class.

\Needspace{13\baselineskip}
\begin{lemma}\label{lem:bounded-closed-potential-nef}
On a connected compact complex manifold $Z$, let $\beta$ be a smooth real $(1,1)$-form and $v$ a bounded quasi-plurisubharmonic function. If $R=\beta+\ddc v\geq0$ is closed, then $\{R\}_{\mathrm{BC}}$ is nef.
\end{lemma}
\begin{proof}
Since $R$ is closed, $\beta$ is closed. Fix a Hermitian metric $\gamma$. By \cref{prop:bounded-almost-regularization}, there are smooth $v_j$ and $\epsilon_j\downarrow0$ such that
\[
 \beta+\ddc v_j\geq-\epsilon_j\gamma.
\]
These are closed representatives of $\{R\}_{\mathrm{BC}}$. Given any $\epsilon>0$, choose $j$ with $\epsilon_j\leq\epsilon$. This is the defining nef estimate.
\end{proof}

\Needspace{8\baselineskip}
\begin{theorem}\label{thm:analytic-current-nef-cartier}\index{Nef class!Cartier realization}
If the closed positive current $T$ has analytic singularities, then $\mathbb D(T)$ is Cartier and has a nef realization on a smooth compact modification.
\end{theorem}
\begin{proof}
Resolve the analytic singularity ideal, using the log-resolution theorem in \cite[section on analytic singularities, theorem on existence of log resolutions]{SGPT}. On one smooth compact model $\pi:Z\to Y$,
\[
 \pi^*T=[D]+R,
\]
where $D$ is a finite effective rational divisor and $R$ is positive, closed, and has locally bounded potentials. To obtain a global bounded potential, choose smooth metrics on the divisor bundles. Subtract the corresponding logarithmic squared norms from $w\circ\pi$. The remainder $v$ is globally bounded, and
\[
 R=\beta_Z+\ddc v
\]
for a smooth closed $\beta_Z$. The lemma makes its class nef.

On every further compact modification $g:Z'\to Z$, the current $g^*R$ still has bounded local potentials and consequently zero Lelong numbers. The entire divisorial part of $g^*\pi^*T$ is therefore $[g^*D]$, and its regular part is $g^*R$. It follows that $\mathbb D(T)_{Z'}=g^*\{R\}_{\mathrm{BC}}$. Thus the pair $(Z,\{R\}_{\mathrm{BC}})$ is the required nef Cartier realization.
\end{proof}

For a general current, analytic regularization gives closed approximants bounded below by $-\epsilon_m\gamma$. Adding $\epsilon_m\gamma$ makes them positive but changes their exterior derivative when $\gamma$ is not closed. The approximation then belongs to the larger spaces of \cref{chap:nonclosed-bclasses}; it does not by itself approximate the original Bott--Chern b-divisor class by nef Cartier classes.


\Needspace{18\baselineskip}
\section{Continuous representatives and pullback}
\label{sec:continuous-bc-representatives}

A finite divisorial sum is pulled back through the local equations of its divisors. For an infinite sum we need pullback to be continuous on closed positive currents whose Bott--Chern classes vary. Solving a scalar equation for the Gauduchon trace selects a smooth representative in each class; the complementary normalized potentials then converge in $L^1$, and \cref{lem:modification-lone-pullback} shows that their pullbacks converge as well.

In this section, $Y$ denotes an arbitrary connected compact complex manifold of dimension $n\geq1$, and $\gamma$ is a Gauduchon metric on $Y$; in dimension one, any Hermitian metric will do. Fix smooth norms on tensor fields induced by $\gamma$. Put $V_\gamma=\int_Y\gamma^n$ and, for a real $(1,1)$-form or current $B$, define its \emph{scalar trace} $\sigma_\gamma(B)$ by
\[
 \sigma_\gamma(B)\,\gamma^n=B\wedge\gamma^{n-1}.
\]
For a current this is an identity of distributions. As in the proof of \cref{thm:gauduchon}, we use the classical elliptic Fredholm theorem, elliptic regularity, and the continuity on distributions of the normalized elliptic inverse.

\Needspace{18\baselineskip}
\begin{lemma}\label{lem:constant-trace-bc-representative}\index{Bott--Chern class!constant-trace representative}
Every closed real $(1,1)$-current $T$ has a unique decomposition
\[
 T=B_\gamma(T)+\ddc v_\gamma(T),\qquad
 \int_Yv_\gamma(T)\,\gamma^n=0,
\]
where $v_\gamma(T)$ is a real distribution and $B_\gamma(T)$ is a smooth closed real form with constant scalar trace. Both assignments are linear and continuous from weak convergence of currents to convergence of distributions. The forms $B_\gamma(T)$ lie in a finite-dimensional space, and their convergence takes place in $C^\infty$.

This space maps isomorphically onto $H^{1,1}_{\mathrm{BC}}(Y,\mathbb R)$. If $T\geq0$, then $v_\gamma(T)$ is quasi-plurisubharmonic. On the cone of positive closed currents, $v_\gamma(T)$ depends continuously on $T$ in $L^1$, and for every nonnegative integer $l$,
\begin{equation}\label{eq:constant-trace-representative-bound}
 \|B_\gamma(T)\|_{C^l}+\|v_\gamma(T)\|_{L^1(\gamma^n)}
       \leq C_l\int_YT\wedge\gamma^{n-1},
\end{equation}
where the constants $C_l$ depend only on the fixed metric and norms.
\end{lemma}
\begin{proof}
Set $L_\gamma f=\sigma_\gamma(\ddc f)$. We use the inverse $G_\gamma$ of \eqref{eq:gaud-trace-green}, normalized with the volume form $\gamma^n/V_\gamma$. It acts on top-degree distributions of total mass zero and is continuous by the elliptic estimates for the operator and its transpose.

\emph{Step 1: construction.} Define
\[
 c_T=\frac{\int_YT\wedge\gamma^{n-1}}{V_\gamma},\qquad
 v_\gamma(T)=G_\gamma\bigl((\sigma_\gamma(T)-c_T)\gamma^n\bigr),\qquad
 B_\gamma(T)=T-\ddc v_\gamma(T).
\]
The current $B_\gamma(T)$ is closed and has scalar trace $c_T$. We now prove its smoothness and the finite-dimensionality assertion without using any K\"ahler identity.

\Needspace{7\baselineskip}
\emph{Step 2: an elliptic system.} Let $J$ be the complex structure acting on real covectors. On real $(1,1)$-forms consider the first-order operator
\[
 \mathscr P B=(\mathrm{d}B,\mathrm{d}\sigma_\gamma(B)).
\]
Its principal symbol at a nonzero real covector $\xi$ is $B\mapsto(\xi\wedge B,\xi\,\sigma_\gamma(B))$, and we claim that it is injective. If $\xi\wedge B=0$, then the $(1,1)$ condition forces $B$ to be a real multiple of $\xi\wedge J\xi$. To check this, choose unitary real coordinates with $\xi=\mathrm{d}x_1$: every term of $B$ contains $\mathrm{d}x_1$, and invariance under $J$ forces every term to contain $\mathrm{d}y_1$ as well. If moreover $\xi\,\sigma_\gamma(B)=0$, then $\sigma_\gamma(B)=0$; since the scalar trace of $\xi\wedge J\xi$ is nonzero, the multiple is zero.

Let $\mathscr P^*$ be the formal adjoint of $\mathscr P$ for the fixed metric and volume. Since the symbol of $\mathscr P$ is injective, $\mathscr P^*\mathscr P$ is elliptic. By elliptic regularity and compactness, its distributional kernel is a finite-dimensional space of smooth forms. Since $\mathscr P B_\gamma(T)=0$, the form $B_\gamma(T)$ belongs to this space. Conversely, the kernel consists exactly of the closed forms with constant scalar trace: integrating $\langle\mathscr P^*\mathscr P B,B\rangle$ gives $\mathscr P B=0$, and connectedness makes the trace constant. Denote this kernel by $\mathscr H_\gamma$.

\emph{Step 3: Bott--Chern representatives.} The construction is linear and continuous on distributions. Convergence of the components $B_\gamma(T)$ therefore takes place in the finite-dimensional space $\mathscr H_\gamma$, hence in every smooth norm. If $T=\ddc f$ for a distribution $f$, then $c_T=0$ and the inverse gives $v_\gamma(T)=f-V_\gamma^{-1}\int_Yf\gamma^n$, so $B_\gamma(T)=0$. Applied to smooth closed forms, this shows that $\mathscr H_\gamma$ contains exactly one representative of each Bott--Chern class. Uniqueness of the decomposition can also be seen directly: if two decompositions are given, $L_\gamma$ of the difference of the two potentials is a constant, integration forces this constant to be zero, and the normalized kernel of $L_\gamma$ is zero.

\emph{Step 4: positive currents.} Suppose now that $T\geq0$. The inequality $\ddc v_\gamma(T)\geq-B_\gamma(T)$ gives a quasi-plurisubharmonic representative of the distribution $v_\gamma(T)$. Let positive closed currents $T_j$ converge weakly to $T$. Their components $B_\gamma(T_j)$ converge smoothly, so the potentials $v_\gamma(T_j)$ have a common smooth lower curvature bound. By \cref{thm:normalized-compactness}, after subtracting their suprema they form an $L^1$-precompact family. If $a_j=\sup_Yv_\gamma(T_j)$, the zero-mean normalization gives
\[
 0\leq a_j V_\gamma=\int_Y\bigl(a_j-v_\gamma(T_j)\bigr)\,\gamma^n,
\]
and the right-hand side is uniformly bounded. Hence the potentials $v_\gamma(T_j)$ themselves are $L^1$-precompact. Every cluster point equals their distributional limit $v_\gamma(T)$, so the whole sequence converges in $L^1$.

Finally, the positive closed currents with trace mass $\int_YT\wedge\gamma^{n-1}=1$ form a weakly compact set: positivity bounds their coefficient measures, weak limits remain positive and closed, and testing with $\gamma^{n-1}$ preserves the trace mass. The smooth and $L^1$ continuity just proved bounds both norms on this set, and positive homogeneity gives \eqref{eq:constant-trace-representative-bound}. A positive current of trace mass zero is zero, so this case is included.
\end{proof}

The smooth representatives give a common curvature bound for converging normalized potentials. This is exactly the hypothesis needed to pull back their $L^1$ convergence.

\Needspace{13\baselineskip}
\begin{proposition}\label{prop:closed-current-modification-continuity}\index{Closed positive current!pullback continuity}
Let $g:Z\to Y$ be a smooth compact modification, and fix a Gauduchon metric $\gamma'$ on $Z$. Pullback of positive closed $(1,1)$-currents, defined through their local potentials, is weakly continuous. Moreover, there is a constant $C_g$, depending only on the modification and the fixed metrics, such that
\begin{equation}\label{eq:modification-linear-mass-bound}
 \int_Zg^*T\wedge(\gamma')^{n-1}
       \leq C_g\int_YT\wedge\gamma^{n-1}
       \qquad(T\geq0,\ \mathrm{d}T=0).
\end{equation}
\end{proposition}
\begin{proof}
Write $T=B_\gamma(T)+\ddc v_\gamma(T)$. The pullback is
\[
 g^*T=g^*B_\gamma(T)+\ddc(v_\gamma(T)\circ g).
\]
This current is positive and closed, since it is locally $\ddc$ of plurisubharmonic functions, and the composed potential is not identically $-\infty$ because $g$ is surjective. The formula agrees with the definition through local potentials, so it does not depend on the decomposition.

If $T_j\to T$, the smooth parts converge in $C^\infty$, and the potentials converge in $L^1(Y)$ with a common lower curvature bound. By \cref{lem:modification-lone-pullback}, their compositions with $g$ converge in $L^1(Z)$, and taking $\ddc$ proves the weak convergence of the pullbacks.

\Needspace{7\baselineskip}
Integration by parts on $Z$ gives
\[
 \int_Zg^*T\wedge(\gamma')^{n-1}
   =\int_Zg^*B_\gamma(T)\wedge(\gamma')^{n-1}.
\]
The absolute value of the right-hand side is at most a fixed constant times $\|B_\gamma(T)\|_{C^0}$. Since the left-hand side is nonnegative, \eqref{eq:constant-trace-representative-bound} gives \eqref{eq:modification-linear-mass-bound}.
\end{proof}

\Needspace{12\baselineskip}
\begin{corollary}\label{cor:countable-divisorial-pullback}\index{Divisorial current!countable pullback}
Let $g:Z\to Y$, $\gamma$ and $\gamma'$ be as in \cref{prop:closed-current-modification-continuity}. Let $D_j$ be prime divisors on $Y$, with integration currents $[D_j]$, and let $c_j>0$. If $S=\sum_j c_j[D_j]$ is a positive closed current of finite mass, then
\[
 g^*S=\sum_j c_j[g^*D_j]
\]
as positive currents of finite mass. The trace masses on $Z$ of the positive tails of this series tend to zero. In particular, $g^*S$ is again purely divisorial.
\end{corollary}
\begin{proof}
Apply the proposition to the increasing finite partial sums $S_N$ and to their positive closed tails $S-S_N$. Pullback is additive on positive closed currents, by addition of local potentials. Thus $g^*S-g^*S_N=g^*(S-S_N)\geq0$, and its trace mass is at most $C_g$ times the trace mass of $S-S_N$, which tends to zero. The pullback of a finite divisor is the usual pullback of its local equations, which identifies $g^*S_N$ with the corresponding partial sum of the series, and hence $g^*S$ with its limit. Its terms are integration currents of prime divisors with nonnegative coefficients, and regrouping this countable positive sum of finite mass does not change it.
\end{proof}


\section{Exceptional divisors between trace models}
\label{sec:trace-exceptional-divisors}

We return to the trace data fixed at the beginning of the chapter, with $\dim Y=r$, and write $S_Z=S_Z(u;Y)$. By \cref{cor:countable-divisorial-pullback}, the full trace current can now be pulled back under a modification $g\colon Z'\to Z$. The pullback is bounded above by the trace current on $Z'$, with equal coefficients along strict transforms, so the difference is a finite effective divisor supported on the exceptional locus. Its logarithmic potential measures the change in the ordinary singularity types of the divisorial parts.

\begin{theorem}\label{thm:trace-exceptional-difference}\index{Exceptional divisor!trace correction and cocycle}
For a smooth compact modification $g\colon Z'\to Z$ over $Y$, there is a unique finite effective real divisor $E_g$ supported on the exceptional locus of $g$ such that
\begin{equation}\label{eq:trace-exceptional-difference}
 S_{Z'}=g^*S_Z+[E_g].
\end{equation}
For a further smooth compact modification $h\colon Z''\to Z'$, these divisors satisfy
\begin{equation}\label{eq:trace-exceptional-cocycle}
 E_{g\circ h}=E_h+h^*E_g.
\end{equation}
Their coefficients are additive and positively homogeneous in the ambient potential $u$.

For every $m$, the analytic-ideal divisors on the two models differ by a finite effective divisor $E_{m,g}$ supported on the exceptional locus:
\[
 [E_{m,g}]=D_{m,Z'}-g^*D_{m,Z}.
\]
\Needspace{5\baselineskip}
For every Gauduchon metric $\gamma'$ on $Z'$, the measures
\[
 [E_{m,g}]\wedge(\gamma')^{r-1}
     \longrightarrow [E_g]\wedge(\gamma')^{r-1}
\]
converge in total variation.
\end{theorem}
\begin{proof}
\emph{Step 1: the exceptional divisor.} Enumerate the prime divisors on $Z$ with positive trace coefficient. By \cref{thm:trace-divisorial-model-compatibility}, the pullback of every finite partial sum of $S_Z$ is at most $S_{Z'}$, and \cref{cor:countable-divisorial-pullback} identifies the increasing limit of these pullbacks with $g^*S_Z$. Thus $S_{Z'}-g^*S_Z$ is a positive closed divisorial current. The two currents have equal coefficients along the strict transform of every prime divisor on $Z$, because these coefficients are computed by the same valuation. Every other prime divisor on $Z'$ is exceptional for $g$, and the exceptional analytic set has only finitely many irreducible components on the compact manifold $Z'$. Hence the difference is the integration current of a finite effective real divisor $E_g$ supported on its divisorial components, and the coefficients of integration currents determine $E_g$ uniquely.

\emph{Step 2: the cocycle identity.} Pullback of positive closed $(1,1)$-currents is functorial, since it is defined by composing local plurisubharmonic potentials. Using \eqref{eq:trace-exceptional-difference} for $h$ and for $g$ gives
\[
 S_{Z''}=h^*S_{Z'}+[E_h]
        =(g\circ h)^*S_Z+[h^*E_g]+[E_h],
\]
and uniqueness gives \eqref{eq:trace-exceptional-cocycle}. Additivity and positive homogeneity follow from the same properties of the trace coefficients and the linearity of pullback.

\Needspace{10\baselineskip}
\emph{Step 3: approximation.} For fixed $m$, apply \eqref{eq:trace-ideal-model-inequality} to the finitely many prime divisors in the support of the image ideal on $Z$. After division by $m$, it gives $D_{m,Z'}-g^*D_{m,Z}\geq0$. The coefficients along strict transforms vanish, so this difference is the integration current of an effective real divisor $E_{m,g}$ supported on the exceptional locus.

Choose also a Gauduchon metric $\gamma$ on $Z$. By \cref{thm:trace-divisorial-current-budget}, the currents $S_Z-D_{m,Z}$ and $S_{Z'}-D_{m,Z'}$ are positive, and their trace masses tend to zero. Subtracting the two identities for the exceptional currents gives
\[
 [E_{m,g}]-[E_g]
   =-(S_{Z'}-D_{m,Z'})+g^*(S_Z-D_{m,Z}),
\]
where both $S_{Z'}-D_{m,Z'}$ and $g^*(S_Z-D_{m,Z})$ are positive. By \eqref{eq:modification-linear-mass-bound}, the total variation of the signed measure $([E_{m,g}]-[E_g])\wedge(\gamma')^{r-1}$ is therefore at most
\[
 \int_{Z'}(S_{Z'}-D_{m,Z'})\wedge(\gamma')^{r-1}
 +C_g\int_Z(S_Z-D_{m,Z})\wedge\gamma^{r-1}.
\]
Both terms tend to zero, which proves convergence in total variation.
\end{proof}

For a modification $g\colon Z'\to Z$ as in the theorem, the coefficient of $E_g$ along an exceptional prime divisor $F\subset Z'$ is
\begin{equation}\label{eq:trace-exceptional-coefficient}
 \tau_{Y,u}(F)
  -\sum_{D\subset Z}\tau_{Y,u}(D)\operatorname{ord}_F(g^*D),
\end{equation}
where $D$ ranges over the prime divisors on $Z$ with positive trace coefficient. This countable sum converges: it is the coefficient of $g^*S_Z$ along $F$ in the positive divisorial series of \cref{cor:countable-divisorial-pullback}, and it is at most $\tau_{Y,u}(F)$. The same current identity compares the ordinary singularity types attached to $S_Z$ and $S_{Z'}$ by \cref{cor:divisorial-current-ordinary-type}.

\Needspace{14\baselineskip}
\begin{corollary}\label{cor:trace-model-ordinary-comparison}\index{Divisorial current!ordinary types under pullback}
Choose decompositions $S_Z=B_Z+\ddc v_Z$ and $S_{Z'}=B_{Z'}+\ddc v_{Z'}$ with smooth closed forms $B_Z,B_{Z'}$, as in \cref{cor:divisorial-current-ordinary-type}. Let $\rho_{E_g}$ be a global logarithmic potential of $E_g$, that is, a quasi-plurisubharmonic function on $Z'$ such that $\ddc\rho_{E_g}-[E_g]$ is smooth. Then
\[
 v_{Z'}-(v_Z\circ g)-\rho_{E_g}
\]
is represented by a smooth function. In particular,
\[
 v_{Z'}\preceq v_Z\circ g
\]
in the sense of \cref{def:singularity-order}, and the two ordinary singularity types are equal if and only if $E_g=0$.
\end{corollary}
\begin{proof}
Global logarithmic potentials exist: choose smooth Hermitian metrics on the line bundles of the prime components of $E_g$, let $\alpha_{E_g}$ be the corresponding weighted sum of their curvature forms, and let $\rho_{E_g}$ be the weighted sum of the logarithms of the squared norms of the canonical sections. Then
\[
 \ddc\rho_{E_g}=[E_g]-\alpha_{E_g},
\]
and any two global logarithmic potentials differ by a smooth function, by \cref{lem:closed-current-global-potential}. We may therefore use this choice. The function $\rho_{E_g}$ is quasi-plurisubharmonic and bounded above on the compact manifold $Z'$.

\Needspace{10\baselineskip}
The chosen decompositions and \eqref{eq:trace-exceptional-difference} give
\[
 \ddc\bigl(v_{Z'}-v_Z\circ g-\rho_{E_g}\bigr)
       =g^*B_Z+\alpha_{E_g}-B_{Z'},
\]
which is smooth. By \cref{lem:closed-current-global-potential}, this difference of locally integrable functions is represented by a smooth function. The resulting equality of quasi-plurisubharmonic functions extends from almost everywhere to everywhere by \cref{lem:qpsh-almost-everywhere-order}. Since $\rho_{E_g}$ is bounded above, this gives the ordinary domination.

If $E_g=0$, the smooth difference gives ordinary equivalence. Conversely, ordinary equivalence would give equal generic Lelong numbers along every prime divisor of $Z'$, whereas the identity just proved shows that along a prime component of $E_g$ these numbers differ by its positive coefficient. Hence ordinary equivalence forces $E_g=0$.
\end{proof}

When the restricted background $\theta|_Y$ is closed, the positive residual currents $R_Z$ constructed in the proof of \cref{prop:trace-bott-chern-bclass}, and their Bott--Chern classes, satisfy
\[
 \begin{split}
 R_{Z'}&=g^*R_Z-[E_g],\\
 \mathbb A_{Z'}(u;Y,\theta)
       &=g^*\mathbb A_Z(u;Y,\theta)-\{[E_g]\}_{\mathrm{BC}}.
 \end{split}
\]
Indeed, the smooth background and the common trace potential pull back, while \eqref{eq:trace-exceptional-difference} gives the change in the subtracted divisorial current. The correction $\{[E_g]\}_{\mathrm{BC}}$ accounts for the additional divisorial vanishing prescribed on the higher model. This is consistent with the exact pullback of the trace envelope in \cref{thm:trace-envelope-modification}, since the envelope imposes the trace inequalities on all models at once. The following example shows that the correction can be nonzero even when the trace data are realized exactly by a potential.

\Needspace{9\baselineskip}
\begin{example}\label{ex:exceptional-trace-with-exact-potential}
Let $Y=\mathbb P^2$ with homogeneous coordinates $[t_0:t_1:t_2]$, let $\sigma$ be the Fubini--Study form normalized by $\int_Y\sigma^2=1$, and choose $a\in\mathbb Q_{>0}$. Put
\[
 u=a\log\frac{|t_1|^2+|t_2|^2}
                   {|t_0|^2+|t_1|^2+|t_2|^2},
\]
and take the ambient manifold to be $Y$ itself. Then $u\in\PSH(Y,a\sigma)$. It has analytic singularities, with its only pole at $p=[1:0:0]$, and zero generic Lelong number on $Y$. By \cref{prop:trace-naive-comparison}, its trace data are given by the generic Lelong numbers of its pullbacks: $\tau_{Y,u}(E)=\nu(u\circ\pi,E)$ for every prime divisor $E$ on a smooth compact modification $\pi$ of $Y$. Since $u$ has no divisorial poles on $Y$, we have $S_Y=0$, whose canonical ordinary type is that of a smooth function. If $\pi\colon\widetilde Y\to Y$ is the blowup at $p$ with exceptional curve $F$, then
\[
 S_{\widetilde Y}=a[F],\qquad E_\pi=aF.
\]
Indeed, on either standard chart of the blowup, the numerator is the squared modulus of the exceptional coordinate times a smooth positive function, so the pole order is $a$ along $F$ and zero along every other prime divisor. Thus the trace data are realized exactly by $u$, although $E_\pi\neq0$. The exceptional difference does not obstruct exact realization; it prevents the identification of the realizing potential with the canonical ordinary type of the divisorial part on the original model, which here is smooth.

\Needspace{8\baselineskip}
The same phenomenon occurs on a manifold without the bounded mass property. Let $(W,\omega_W)$ be a connected compact Hermitian manifold of dimension $k$ that fails the bounded mass property, such as the Hopf threefold of \cref{thm:hopf-unbounded}, and pull $u$ back to $Y\times W$; take both the ambient manifold and the submanifold to be $Y\times W$. The pole set of the pullback is $\{p\}\times W$, of codimension two, and the modification $\widetilde Y\times W\to Y\times W$ has exceptional coefficient $a$ along $F\times W$. Since the pullback of $u$ has analytic singularities, \cref{prop:trace-naive-comparison} again shows that it realizes the trace data exactly. Let $\operatorname{pr}_Y$ and $\operatorname{pr}_W$ be the two product projections.

\Needspace{9\baselineskip}
Choose smooth $\omega_W$-plurisubharmonic functions $H_j$ whose masses tend to infinity. For the Hermitian form $\Omega=\operatorname{pr}_Y^*\sigma+\operatorname{pr}_W^*\omega_W$, the pullbacks of the $H_j$ have masses
\[
 \binom{k+2}{2}\left(\int_Y\sigma^2\right)
          \left(\int_W(\omega_W+\ddc H_j)^k\right)\longrightarrow+\infty,
\]
by the binomial expansion of smooth forms and the dimensions of the two factors. Hence the product also fails the bounded mass property. The pulled-back potential belongs to $\PSH(Y\times W,C\Omega)$ for every $C\geq\max\{a,1\}$, since its curvature is bounded below by $-a\operatorname{pr}_Y^*\sigma$. Scaling the smooth functions $H_j$ by $C$ shows that $C\Omega$ also has unbounded mass.
\end{example}

